\section{Continuum substrate model}
\label{sec:continuum-model}

\subsection{Kinematics and induced metric}

Let $x=(x^1,x^2,x^3)\in\Omega\subset\R^3$ denote reference coordinates of the brane.
The embedding $\vecX(x,t)\in\R^4$ induces a (time-dependent) metric on the brane:
\begin{equation}
  g_{ij}(x,t) := \partial_i \vecX(x,t)\cdot \partial_j \vecX(x,t), \qquad i,j=1,2,3.
  \label{eq:induced-metric}
\end{equation}
Rotational invariance within the brane follows if the stored energy depends on invariants of $g$
relative to a reference metric $\bar g$.

\paragraph{Pre-tension as a reference-metric mismatch.}
To model a homogeneous pre-tension (the continuum analogue of a prestretched lattice),
we introduce a reference metric $\bar g_{ij}=\alpha^2\,\delta_{ij}$ with $\alpha\in(0,1]$.
Embedding the brane in a configuration close to $g_{ij}\approx \delta_{ij}$ then represents an
initial state that is ``larger'' than its preferred metric and therefore carries tension.
This parametrization is an estimator-level device; alternative but equivalent parametrizations
(e.g., multiplicative pre-deformation gradients) can be substituted \citep{Ogden1984,Holzapfel2000}.

\subsection{Action principle}

We consider an action
\begin{equation}
  S[\vecX] = \int \mathrm{d}t \int_{\Omega} \mathrm{d}^3x\;
  \mathcal{L}\big(\vecX,\partial_t \vecX,\partial_i \vecX\big),
  \qquad
  \mathcal{L} := \frac{\rho_m}{2}\,|\partial_t \vecX|^2 - W(g;\bar g),
  \label{eq:action}
\end{equation}
where $\rho_m$ is a mass density and $W$ is an isotropic stored-energy density.
A minimal requirement is that $W$ is frame-indifferent and stable around the background,
so that a linear wave regime exists.

\paragraph{Example constitutive families.}
The analysis is compatible with standard isotropic hyperelastic families.
For instance, $W$ may be expressed via invariants of the right Cauchy--Green tensor
(relative to $\bar g$), and reduces to linear elasticity for small strains \citep{Ogden1984,Holzapfel2000}.
We keep the constitutive choice modular; the diagnostic pipeline later depends mainly
on the existence of polarized, narrowband propagating modes.

\subsection{Equations of motion}

Define the first Piola--Kirchhoff-like stress
\begin{equation}
  P_{Ai} := \frac{\partial W}{\partial(\partial_i X^A)} ,\qquad A=1,\dots,4,
  \label{eq:piola}
\end{equation}
so that the Euler--Lagrange equations take the form
\begin{equation}
  \rho_m\,\partial_t^2 X^A = \partial_i P_{Ai}.
  \label{eq:eom}
\end{equation}
Because $W$ is constructed from $g_{ij}$ (and $\bar g_{ij}$), the continuum model is rotationally
invariant in the brane coordinates by construction.

\subsection{Linearization and polarized modes}

Linearize around a flat background $\vecX_\star(x)=(x,0)$ and small displacements
$\vecu(x,t)=\vecX(x,t)-\vecX_\star(x)$. Under standard regularity assumptions on $W$,
one obtains a linear wave system with (in general) multiple branches, including transverse and
longitudinal polarizations. These are the substrate carriers for the later geometric-phase diagnostics.
